\documentclass[12pt, uplatex, dvipdfmx, article]{jlreq}
\usepackage{amssymb,amsmath,latexsym}
\usepackage[skins]{tcolorbox}
% \usepackage[haranoaji]{pxchfon} % 原ノ味フォントを使う


% 5つの引数を受け取る \present コマンドの定義
\newcommand{\speak}[5]{ %内部的にtcolorboxを使っている
\begin{tcolorbox}[
    colback= white,      % 背景色を白
    colframe= white,     % 枠線を白色に
    colbacktitle= white,
    before skip= 1pt,    % ボックス外側の「上」余白
    after skip= 1pt,     % ボックス外側の「下」余白
    top= 0.7mm,            % ボックス内側の「上」余白
    bottom= 0.7mm,         % ボックス内側の「下」余白
    sidebyside,               % ボックスを左右に分割
    sidebyside align= top,     % 上端を揃える
    lefthand width= 28mm,      % 左側（時間）の幅
    sidebyside gap= 3mm,       % 時間と講演者の間の隙間
    lower separated= false     % 中央の区切り線を消す
    ]
    {#1}\hspace{1.5mm}～\hspace{1.5mm}{#2} %#1：開始時間、#2：終了時間
    \tcblower                 % ここから右側（講演者・タイトル）
    {#3}(#4)\\ %#3：講演者名、#4：所属・学年
    {#5} %5：講演タイトル
\end{tcolorbox}
}


\begin{document}
\pagestyle{empty} %ページ番号を消す

\begin{center}
{\huge %研究集会名
    ConferenceName
}
\vspace{12pt}

%他詳細
\begin{tabular}{l}
    {日程：yyyy年mm月dd日 hh:mm～hh:mm}\\
    {会場：where}\\
    {形式：style}\\
    {Zoomリンク：}
\end{tabular}

\vspace{10pt}
    プログラム
\end{center}

% \speak{開始時間}{終了時間}{講演者}{所属}{学年}{講演タイトル}

%Sampleデータなので実際に出力する際には消す
\speak{12:00}{13:00}{山田太郎}{A大学・M1}{微分方程式について}
\speak{13:00}{14:00}{佐藤花子}{B大学・B4}{整数論とゼータ関数の関係}


%これより上に講演内容を書く

\begin{flushleft} %世話人の名前(所属・学年)
    世話人：XXX(YYY・M1)
\end{flushleft}
\end{document}